\section{General Equilibrium: AS-AD}

\begin{frame}{AS-AD Model}
    \mycomment{
    The IS-LM model has given us a foundation to derive the aggregate demand function $AD$, which consists of all the possible intersections of $IS$ and $LM$ at different price levels.
    }
    \showmytikz[1]{../shared/tikz/tikz-asad}
\end{frame}


\begin{frame}{AS-AD Model}
    \mycomment{
    This must be met by an aggregate supply curve. This is far more difficult to derive. 
    
    In the short run, with sticky factor prices and given stock of machinery and capital, supply is inelastic. 

    But supply might be elastic in some situations and in the medium run; particularly, when some effect has left capacity unused, or when it can be satisfied by imports.

    In the long run, aggregate supply is given by the maximum output an economy can achieve, that is the technology level and the size of the labour force, which is typically inelastic.
    }
    \showmytikz[1,3,4.5,7]{../shared/tikz/tikz-asad}
\end{frame}

% \begin{frame}{AS-AD Model}
%     \mycomment{
%     Whenever the economy is a disequilibrium, it will revert to that equilibrium of general price level and BIP - in the short run.
%     }
%     \showmytikz[1,2,3,4,5]{../shared/tikz/tikz-asad}
% \end{frame}

\begin{frame}{AS-AD Model}
    \mycomment{
    If some government policy - fiscal policy, or monetary policy - shifts the equilibrium possibilities (as derived in the IS curve), the aggregate demand curve shifts right.

    What happens now depends on the elasticity of supply. If supply is inelastic, the main effect is an increase in the price level (inflation).
    }
    \showmytikz[1,3,5,6]{../shared/tikz/tikz-asad}
\end{frame}

\begin{frame}{AS-AD Model}
    \mycomment{
    If, however, supply is elastic, the increase in government spending or the easy money policy of the central bank can result in an increase in BIP, while the price level changes only a little.
    }
    \showmytikz[1,3,5,6,7,8]{../shared/tikz/tikz-asad}
\end{frame}
